%% =============================================================================
%% SKELETON - Section 6: Cross-Strategy Interdependencies
%% =============================================================================
%% Tables/figures INLINE (working draft for supervisor).
%% Same structure as paper1_sec6_interdependencies.tex.
%% Maps to: rq3_01 through rq3_07
%% =============================================================================


%% ---- 6.1 Correlations, Co-Widening, and Persistence ----
\subsection{Correlations, Co-Widening, and Persistence}
\label{sec:rq3_corr}

This section examines unconditional and regime-dependent co-movement across
the three arbitrage strategies, using mispricing magnitudes $m_{i,t}$ (basis levels)
rather than returns, following \citet{duffie2010presidential}.

\paragraph{Mispricing Construction.}
We construct $m_{i,t}$ as the absolute basis level for each strategy, sign-adjusted
so that $\Delta m > 0$ denotes widening (increased mispricing) for all three.
BTP~Italia: cross-sectional mean with maturity filter. CDS-Bond: top negative-basis
names. iTraxx Skew: mean of on-the-run series. Regime definition uses iTraxx Main
5Y with fixed thresholds (see Section~3).

\paragraph{Unconditional and Regime-Dependent Correlations.}
Table~\ref{tab:rq3_unified_corr} reports Pearson correlations for both
strategy returns (Panel~A) and mispricing changes $\Delta m$ (Panel~B),
unconditionally and across three stress regimes (LOW/MEDIUM/HIGH,
iTraxx Main 5Y thresholds at 60 and 120 bps).
Correlations display clear monotonicity: the corporate credit pair
(CDS--Bond/iTraxx) becomes more positively correlated as stress
increases, while pairs involving BTP~Italia become progressively more
negative. This two-cluster pattern is interpreted in Section~\ref{sec:rq3_discussion}.
Figure~\ref{fig:rq3_regime_bar} visualizes the regime contrast.

\input{\tablepath/RQ3_Unified_Correlations_article}

\begin{figure}[H]
\centering
\includegraphics[width=0.9\textwidth]{../results/rq3_duffie/figures/C1_regime_correlations_bar.pdf}
\caption{Pairwise Correlations by Stress Regime.
  Bars show Pearson correlations of $\Delta m$ in LOW, MEDIUM, and HIGH
  stress regimes defined by iTraxx Main 5Y thresholds (60/120 bps).}
\label{fig:rq3_regime_bar}
\end{figure}

%% ---- 6.2 Spanning Regressions Across Strategies ----
\subsection{Spanning Regressions Across Strategies}
\label{sec:rq3_crossreg}

\paragraph{Cross-Strategy Spanning.}
Following \citet{fleckenstein2014tips}, we estimate spanning
regressions of the form
\begin{equation}\label{eq:spanning}
  \Delta m_{i,t} = \alpha
  + \sum_{j \neq i} \sum_{l=0}^{L} \beta_{j,l}\, \Delta m_{j,t-l}
  + \varepsilon_t,
\end{equation}
where $\Delta m_{i,t}$ is the monthly change in mispricing for strategy
$i$ and $L = 2$ lags. Inference uses Newey--West HAC standard errors.
Table~\ref{tab:rq3_spanning} reports the results. Significant
contemporaneous coefficients indicate that mispricing changes co-move
beyond what common macro factors would predict. Significant lagged
coefficients point to lead-lag transmission across strategies,
consistent with slow capital reallocation: the more liquid market
(iTraxx index) adjusts first, and the effect propagates to the less
liquid markets (CDS--Bond, BTP~Italia) with a delay of one to two
months. The $\bar{R}^2$ is highest for BTP~Italia, indicating that
the corporate credit cluster has meaningful predictive power for the
inflation-linked basis despite the two-cluster segmentation.
\input{\tablepath/RQ3_Spanning_Regressions_article}

\paragraph{Stress Amplification.}
\citet{duffie2010presidential} predicts that cross-strategy transmission
should intensify when intermediary capital is under pressure.
Table~\ref{tab:rq3_interaction} tests this via interaction regressions:
\begin{equation}\label{eq:interaction}
  \Delta m_{i,t} = \beta_0\, \Delta m_{j,t}
  + \beta_1 (\Delta m_{j,t} \times z_t)
  + \gamma\, z_t + \varepsilon_t,
\end{equation}
where $z_t$ is the standardized iTraxx Main 5Y spread, proxying for
funding stress. The coefficient $\beta_1$ captures the \emph{incremental}
transmission during stress: $\beta_1 > 0$ means that a given shock to
$\Delta m_j$ has a larger effect on $\Delta m_i$ when intermediary
balance sheets are strained. Results confirm the prediction for the
corporate credit cluster, while pairs involving BTP~Italia exhibit
significant \emph{negative} amplification---stress widens the gap between
the institutional and retail segments, consistent with the two-cluster
structure documented above.

\input{\tablepath/RQ3_Interaction_Regressions_article}

\paragraph{Common Factor Structure.}
Table~\ref{tab:rq3_pc1_funding} reports the regression of PC1 (first
principal component of the three $\Delta m$ series) on intermediary
stress proxies, following the logic of \citet{duffie2010presidential}:
if the common factor driving correlated mispricings reflects
intermediary balance-sheet constraints, it should be explained by
variables that capture dealer capital, funding costs, and market
illiquidity. Panel~A selects one proxy per channel from the
intermediary asset pricing literature:
the \citet{he2017intermediary} capital ratio (HKM\_IC),
the Libor--OIS spread as a measure of global dealer funding cost,
European prime broker CDS spreads as a direct gauge of dealer health,
and the \citet{amihud2015pricing} illiquidity measure.
The majority of proxies are significant with signs consistent with
theory: a decline in intermediary capital, a rise in funding costs,
or an increase in market illiquidity all predict wider mispricings.
Panel~B confirms these findings with an extended set of proxies,
providing direct evidence for the slow-moving capital channel.


\input{\tablepath/RQ3_PC1_Funding_article}


%% ---- 6.3 VAR Analysis and Impulse Responses ----
\subsection{VAR Analysis and Impulse Responses}
\label{sec:rq3_var}

\paragraph{Granger Causality.}
The results reveal a clear transmission hierarchy: the iTraxx index
--- the most liquid instrument in the system --- Granger-causes
BTP~Italia, and CDS--Bond Basis also Granger-causes BTP~Italia.
This is consistent with a liquidity pecking order in which
funding shocks first affect the standardised, exchange-traded credit
index, then propagate to the less liquid cash-CDS basis and finally
to the retail-dominated inflation-linked market.

\paragraph{Generalized Impulse Responses.}
Figure~\ref{fig:rq3_girf} plots Generalized IRFs
\citep{pesaran1998generalized}, which are invariant to variable ordering.
Cross-responses confirm the Granger causality findings: a shock to
iTraxx $\Delta m$ generates a significant negative response in BTP~Italia
that persists for approximately two to three months, consistent with the
two-cluster structure. Own-responses on the diagonal reveal that
CDS--Bond shocks are the most persistent (half-life of approximately
four months), while iTraxx shocks dissipate within two months ---
consistent with the liquidity hierarchy.

\begin{figure}[H]
\centering
\includegraphics[width=\textwidth]{../results/rq3_duffie/figures/E1_girf.pdf}
\caption{Generalized Impulse Response Functions.
  Responses of $\Delta m$ to one-standard-deviation shocks.
  Shaded areas: 95\% bootstrap confidence intervals (1000 replications).}
\label{fig:rq3_girf}
\end{figure}



%% ---- 6.4 Discussion: Evidence for Slow-Moving Capital ----
\subsection{Discussion: Evidence for Slow-Moving Capital}
\label{sec:rq3_discussion}

\paragraph{Investor Base Segmentation: The Two-Cluster Structure.}
CDS--Bond and iTraxx share the same institutional participants (hedge funds,
dealer desks, leveraged accounts) and co-widen in stress as common funding
shocks force simultaneous unwinding \citep{brunnermeier2009market}.
BTP~Italia is a retail instrument with a loyalty bonus
(\textit{premio fedelt\`{a}}) that rewards hold-to-maturity behaviour,
explaining why retail investors do not sell even when the basis offers
apparent value relative to nominal BTPs.
These preferred-habitat investors \citep{vayanos2021preferred} do not use
leverage or face margin calls, so the BTP~Italia basis does not participate
in institutional fire sales---explaining the negative correlation
with the corporate credit cluster that intensifies monotonically
from LOW to HIGH stress.
 
\paragraph{Why Do These Arbitrages Persist?}
Each strategy reflects a distinct limit to arbitrage: retail investor
segmentation for BTP~Italia \citep{vayanos2021preferred}; funding risk,
counterparty risk, and bond illiquidity for CDS--Bond
\citep{bai2019cds, mitchell2012arbitrage}; and post-crisis
regulatory capital costs for the iTraxx skew, where the SLR reduces
dealer ROE to $\approx$3\%
\citep{boyarchenko2016trends, oehmke2017anatomy}.
The persistence of alpha documented in Sections~4--5 is consistent with
slow-moving capital keeping mispricings open longer than standard
models predict \citep{duffie2010presidential, mitchell2007slow}.